\section*{局限性}\label{sec:limitations}

主要实验的学生模型较小，三个学生种子对初始化变异的覆盖有限。每个领域内的语料与教师模型固定；学生模型的重复实验并非教师模型的重复实验。对照共享 CE 与 TG 最终结果，并非相互独立的析因实验。尚无完整析因干预能够分离教师质量、架构、错误结构与优化的影响。教师状态干预同时改变了教师训练历程与预测质量。TT/TG 在似然、参数量及训练方面的不匹配，使其无法成为单独考察架构影响的对照。

裁剪警告及非零的 AMP 跳步比例限制了基于优化的解释。仅使用 T 的监督在训练早期每次更新都触发裁剪，现代模型实验的 KD5 始终处于完全裁剪饱和状态。本研究没有评估自适应 KD 基线，因此不能确立相对于自适应方法的优势。局部诊断指标集合规模较小、存在依赖关系且采用回顾性分析，而非前瞻性的教师选择基准。

现代模型压力测试处于退化的续训设置中，仅使用 CE 的续训改进空间门控也未通过。SmolLM2 的预训练包含源自 FineWeb 的数据，但与所选子集在文档层面的确切重叠仍未知。现代模型的各个种子分别独立适配同一个预训练学生基座，而非重复预训练过程。结论能否推广到更大的学生模型、其他教师模型、生成任务指标及不同续训设置，仍有待研究。
